\documentclass[11pt,a4paper]{article}
\usepackage{amsmath,amssymb,amsthm}
\usepackage[margin=2.5cm]{geometry}
\usepackage{hyperref}

\newtheorem{definition}{Definition}[section]
\newtheorem{theorem}[definition]{Theorem}
\newtheorem{proposition}[definition]{Proposition}
\newtheorem{lemma}[definition]{Lemma}
\newtheorem{corollary}[definition]{Corollary}
\newtheorem{conjecture}[definition]{Conjecture}
\newtheorem{remark}[definition]{Remark}

\title{Hyperseed Formalization:\\
  The Anthropic Fable Farce}
\author{ProtomegaTron (automated formalization)}
\date{2026-09-17}

\begin{document}
\maketitle

\begin{abstract}
We formalize the intellectual structure of Ben Goertzel's Substack article
``The Anthropic Fable Farce'' (2026-07-08) within the Hyperseed ontology.
The article analyzes the 48-hour episode in which Anthropic's Fable~5 was
shipped, export-controlled, and pulled as a controlled demonstration that
centralized AI control fails on contact with real-world geopolitics. Each
structural insight maps onto Hyperseed structures: state coercion as
external cocycle override, KYC as provenance theatre, capability diffusion
as base space equalization, arms-race thinking as zero-sum fiber collapse,
and the distinction between structural (fiber-level) and institutional
(section-level) safety.
\end{abstract}

\tableofcontents

%% =================================================================
\section{The Fable episode as controlled experiment}
\label{sec:experiment}
%% =================================================================

\begin{definition}[External cocycle override --- claims 1--6, 22]
\label{def:override}
The Fable episode demonstrates a fundamental vulnerability of
centralized AI systems: a single government directive can override
the provider's own policies and kill global access.

In Hyperseed terms, this is an \emph{external cocycle override}: the
transition functions of the capability bundle are rewritten by an
authority external to the bundle's own governance structure:
\[
  g_{ij}^{\text{Anthropic}} \xrightarrow{\text{directive}}
  g_{ij}^{\text{government}}
\]

The centralized architecture makes this possible because the entire
bundle passes through a single control point (Anthropic's API). Any
entity that can coerce that control point can rewrite the bundle's
transition functions. The override works regardless of the provider's
intentions --- Anthropic held principled red lines (no mass surveillance,
no autonomous weapons) but was overridden anyway.

The episode demonstrates three structural vulnerabilities of
centralized control:
\begin{enumerate}
\item \textbf{Single point of coercion.} All users depend on one API;
  controlling that API controls all users.
\item \textbf{Principled refusal is insufficient.} The provider's
  principles can be overridden by controlling its market access,
  regulatory status, or supply chain designation.
\item \textbf{Collateral damage is total.} The directive intended to
  restrict foreign nationals ended up restricting everyone --- the
  granularity of centralized control is too coarse for targeted
  policy.
\end{enumerate}

This is why structural (fiber-level) safety is needed: it doesn't
depend on any single institution maintaining good behavior under
political pressure.
\end{definition}

%% =================================================================
\section{KYC as provenance theatre}
\label{sec:kyc}
%% =================================================================

\begin{proposition}[Broken provenance --- claims 10--13, 23]
\label{prop:kyc}
If the export-control policy sticks, the endpoint is KYC (Know Your
Customer) for frontier AI: upload your passport to use the latest
model. The article argues this is fundamentally broken:

\begin{enumerate}
\item \textbf{Leaky barrier.} Criminals buy verified accounts ---
  a staple of money laundering for decades, now enhanced by AI.
\item \textbf{Deepfake liveness.} Liveness checks (hold up passport,
  turn head) are defeated by deepfakes that paste another person's
  face convincingly.
\item \textbf{Honeypot creation.} Running KYC requires warehousing
  passports and biometrics --- creating a target that, when breached,
  feeds the identity-forgery economy.
\item \textbf{Wrong target.} KYC stops casual users but not
  determined adversaries (state actors, terrorist networks) who
  defeat identity controls at scale.
\end{enumerate}

In Hyperseed terms, KYC is \emph{provenance theatre}: it produces
provenance cues (verified documents, liveness checks) that the system
treats as establishing identity, but these cues are systematically
defeatable. The provenance chain \emph{looks} clean --- the system
records a verified identity --- but the actual identity behind the
verification may be forged.

This is the same $R_{\text{auth}}$ failure mode as deceptive open
letters: the stated provenance (verified identity) doesn't match the
actual provenance (purchased account, deepfaked face). The failure is
structural, not accidental --- the system is designed to accept
provenance cues that are known to be defeatable.

A genuine provenance system would need to verify identity through
means that are harder to forge: cryptographic keys tied to actual
capabilities (proof of asset control), reputation histories that
accumulate over time and cannot be instantly purchased, or zero-knowledge
proofs of properties without revealing identity. These are the
provenance mechanisms being developed in the decentralized AI ecosystem.
\end{proposition}

%% =================================================================
\section{Capability diffusion as base space equalization}
\label{sec:diffusion}
%% =================================================================

\begin{theorem}[Structural unsustainability --- claims 14--16, 24]
\label{thm:diffusion}
The article argues that capability containment through monopoly is
structurally unsustainable: whatever the largest model can do, smaller
ones manage months later; open Chinese models achieve parity; the
ban shifts market share without denying adversaries anything.

In Hyperseed terms, this is \emph{base space equalization}: the base
space of model capability equalizes across providers over time through
multiple mechanisms:
\begin{enumerate}
\item \textbf{Scaling law convergence.} Given sufficient compute,
  different architectures converge to similar capability profiles.
\item \textbf{Knowledge diffusion.} Techniques (training methods,
  data curation, architecture innovations) spread through papers,
  reverse engineering, and employee movement.
\item \textbf{Open model ecosystem.} Open-weight releases from
  multiple countries (especially China) provide baseline capability
  that any actor can build on.
\end{enumerate}

The equalization dynamic means that the base space $\mathcal{B}$
(capability) becomes roughly uniform across providers:
\[
  \lim_{t \to \infty} \max_{i,j}
  d(\mathcal{B}_i(t), \mathcal{B}_j(t)) \to 0
\]
where the convergence timescale is months, not years.

Containment through capability monopoly requires maintaining
$d(\mathcal{B}_{\text{leader}}, \mathcal{B}_{\text{rest}}) > \epsilon$
indefinitely. The equalization dynamic makes this structurally
unsustainable: you can delay equalization (through export controls,
classification, hiring restrictions) but not prevent it, because the
underlying knowledge is too distributed and the replication cost
decreases with each generation.

The practical consequence: the export-control directive doesn't deny
adversaries the capability; it shifts where the capability is
developed and who controls it. ``Three or six months from now, China
releases a Mythos++ and hands it to everyone.''
\end{theorem}

%% =================================================================
\section{Arms-race thinking as zero-sum fiber collapse}
\label{sec:arms}
%% =================================================================

\begin{proposition}[Scalar collapse in geopolitics --- claims 17--19, 25]
\label{prop:arms}
The article identifies arms-race thinking (``we must keep the edge
from our rivals'') as the deeper mistake underlying the entire episode.

In Hyperseed terms, arms-race thinking is \emph{zero-sum fiber
collapse}: reducing the rich multi-fiber structure of AI development
(many independent development paths, each carrying different values,
architectures, governance structures, and safety approaches) to a
single competitive dimension (``who has the most powerful model'').

This is scalar collapse applied to geopolitics:
\[
  \underbrace{(v_1, v_2, \ldots, v_n)}_{\text{multi-fiber: values,
  architectures, governance}} \xrightarrow{\text{arms-race frame}}
  \underbrace{s}_{\text{scalar: who's winning}}
\]

The collapse loses exactly the structure that makes safety possible:
\begin{itemize}
\item \textbf{Diversity of approaches.} Different safety strategies
  (neural-symbolic verification, constitutional AI, RLHF, formal
  methods) are all flattened to ``does it help us win?''
\item \textbf{Governance variation.} Different governance structures
  (open, corporate, academic, decentralized) are all evaluated on
  competitive contribution, not safety contribution.
\item \textbf{Value pluralism.} Different value systems and cultural
  perspectives on AI development are collapsed to national-interest
  alignment.
\end{itemize}

The positive-sum alternative preserves the fiber structure: build
systems that are safe by construction (structural safety), open
enough that safety is verifiable (provenance transparency), and
distributed enough that no single directive can switch them off
(distributed governance). This frame treats the fiber structure as
a feature, not a liability.

The security arms race (models getting better at both attacking and
defending) is best tamed by \emph{correct-by-construction software}
--- moving the base space itself to a domain where the attack surface
is structurally reduced, rather than trying to stay ahead in the
arms race within the current base space.
\end{proposition}

%% =================================================================
\section{Structural vs.\ institutional safety}
\label{sec:structural}
%% =================================================================

\begin{definition}[Fiber vs.\ section --- claims 20--21, 27]
\label{def:structural}
The article's central thesis: ``open and decentralized is the only
kind of safety that survives contact with the real world.''

In Hyperseed terms, this is the distinction between \emph{fiber-level}
and \emph{section-level} safety:

\begin{itemize}
\item \textbf{Structural (fiber-level) safety.} Built into the
  system's architecture: verification mechanisms, distributed
  governance, open-weight inspectability, cryptographic provenance.
  This safety is preserved under any base-space transformation ---
  change the government, the company, the market conditions, and
  the safety mechanisms still function because they're in the fiber,
  not in any particular section.

\item \textbf{Institutional (section-level) safety.} Dependent on
  particular institutions maintaining good behavior: Anthropic's
  red lines, a government's regulations, a CEO's principles. This
  safety lives in particular sections of the bundle --- it exists
  at specific base points (specific political configurations,
  specific corporate leadership) and can be overridden by changing
  the base point (new administration, new directive, acquisition,
  coercion).
\end{itemize}

The Fable episode is a demonstration of section-level safety failing:
Anthropic's principled policies (a section) were overridden by a
government directive (base-point change). Fiber-level safety would
have been unaffected --- if the safety mechanisms are in the
architecture rather than the policy, no directive can remove them
without rebuilding the system.

This maps onto a general Hyperseed principle:
\begin{quote}
Any safety property that depends on a particular entity maintaining
good behavior is section-level safety and will eventually be
overridden. Only fiber-level safety --- safety built into the
mathematical and architectural structure of the system --- survives
the full range of base-space transformations that the real world
will subject it to.
\end{quote}
\end{definition}

%% =================================================================
\section{Danger theatre and cocycle contradiction}
\label{sec:theatre}
%% =================================================================

\begin{proposition}[Own-goal dynamics --- claims 7--9, 26]
\label{prop:theatre}
The article identifies a cocycle contradiction in Anthropic's strategy:
\begin{enumerate}
\item Anthropic promoted a transition function $g_1$: ``our models are
  world-historically dangerous'' (danger framing for marketing and
  policy influence).
\item The government applied its own transition function $g_2$:
  ``if it's that dangerous, we'll export-control it'' (national
  security response to stated danger).
\item The composition $g_2 \circ g_1$ produces an outcome that
  contradicts Anthropic's intended trajectory: global access to
  their flagship is killed.
\end{enumerate}

The cocycle doesn't close the way Anthropic expected. They intended
$g_1$ to establish their brand as powerful-but-responsible; instead,
the government took the ``dangerous'' part at face value and acted
on it. The transition functions composed in a direction Anthropic
didn't control.

This is a general warning about danger framing: once you inject
the transition function ``this is too dangerous for unrestricted
access'' into the political cocycle, you cannot control how other
actors compose with it. The function is public and can be applied
by anyone with the authority to restrict access --- and the
government has more authority than the company.
\end{proposition}

%% =================================================================
\section{Cross-references}
%% =================================================================

\begin{remark}[Connections to other formalizations]
\begin{itemize}
\item \textbf{Fable of Benevolent Throttling} (2026-07-12): same
  episode analyzed at a deeper structural level. The throttling
  analysis identifies the cocycle ambiguity (safety vs.\ competition);
  this article focuses on the geopolitical dynamics.
\item \textbf{Un-Jailbreakable AI} (2026-07-10): constructive
  alternative. Where this article diagnoses the failure of centralized
  control, the un-jailbreakable article provides the neural-symbolic
  generate-and-verify architecture as the constructive response.
\item \textbf{Geopolitics of the Great AI Bet} (2026-08-13):
  geopolitical context. The arms-race frame and government-corporate
  fusion analyzed here are part of the broader geopolitical dynamics.
\item \textbf{Why I Didn't Sign} (2026-07-15): provenance theatre.
  KYC as broken provenance is the same structural pattern as
  deceptive open letters --- stated provenance doesn't match actual
  provenance.
\item \textbf{Proof of Humanity} (2026-08-04): alternative identity
  verification. The cryptographic provenance mechanisms described
  there (zero-knowledge proofs, reputation histories) are the
  constructive alternative to KYC's broken provenance.
\item \textbf{Big Tech's Big Bets} (2026-07-25): competitive
  dynamics. The arms-race frame and danger-as-marketing pattern
  are part of the broader big-tech competitive landscape.
\item \textbf{Decentralized Digital} (2026-08-06): structural remedy.
  The open/decentralized infrastructure being built (ASI Alliance,
  SingularityNET) is the fiber-level safety alternative to
  institutional (section-level) safety.
\end{itemize}
\end{remark}

\begin{conjecture}[Fiber-level safety sufficiency]
Conjecture: a system with fiber-level safety (safety mechanisms
built into architecture, not dependent on institutional behavior)
is robust to any base-space transformation that preserves the fiber
structure --- i.e., any change that doesn't require rebuilding the
system from scratch. Formally:
\[
  \forall \phi : B \to B' \text{ (base change)}, \quad
  \text{Safety}(E, \nabla) \implies
  \text{Safety}(\phi^* E, \phi^* \nabla)
\]
where $\phi^* E$ is the pullback bundle and $\phi^* \nabla$ the
pullback connection. The safety property is preserved because it
lives in the fiber, not at any particular base point.

Institutional (section-level) safety fails this test: it's a property
of specific sections $s : B \to E$, and there's no guarantee that
$\phi^* s$ exists or preserves the safety property.
\end{conjecture}

\end{document}
